\FloatBarrier
\subsection{Sensitivity to missing-data handling}
The conclusions were robust to the imputation strategy (Table~\ref{tab:missing}). With iterative or $k$-NN imputation, ROC-AUC changed by at most $+0.005$ relative to median imputation for both logistic regression and RFECV + stacking (all $p \geq 0.45$). Adding missingness indicators did not help. It lowered the ROC-AUC of logistic regression slightly ($-0.010$, $p = 0.39$) and left RFECV + stacking unchanged ($-0.002$, $p = 0.50$), and RFECV retained on average only 1.2 of the five indicators. Under every strategy, stacking and logistic regression remained within 0.005 of each other.

\begin{table*}[!htbp]
\centering
\caption{Missing-data sensitivity analysis on the development set (10-fold cross-validation repeated twice; 20 folds). ROC-AUC values are means with corrected 95\% CIs. $\Delta$AUC is the paired difference from median imputation for the same model, with the corrected resampled $t$-test $p$-value.}
\label{tab:missing}
\resizebox{\textwidth}{!}{%
\begin{tabular}{llcccccl}
\toprule
Model & Imputation & Features & ROC-AUC (95\% CI) & MCC & Sens. & Brier & $\Delta$AUC ($p$) \\
\midrule
Logistic regression & Median & 16.0 & 0.846 (0.813--0.878) & 0.514 & 0.743 & 0.162 & reference \\
Logistic regression & Median + indicators & 21.0 & 0.836 (0.788--0.884) & 0.484 & 0.727 & 0.166 & -0.010 (0.387) \\
Logistic regression & Iterative (MICE-style) & 16.0 & 0.847 (0.814--0.879) & 0.517 & 0.753 & 0.162 & +0.001 (0.712) \\
Logistic regression & KNN ($k=10$) & 16.0 & 0.847 (0.813--0.880) & 0.505 & 0.743 & 0.162 & +0.001 (0.607) \\
\midrule
Stacking + RFECV & Median & 11.3 & 0.844 (0.804--0.883) & 0.498 & 0.752 & 0.164 & reference \\
Stacking + RFECV & Median + indicators & 13.3 & 0.841 (0.799--0.884) & 0.496 & 0.748 & 0.166 & -0.002 (0.503) \\
Stacking + RFECV & Iterative (MICE-style) & 14.0 & 0.848 (0.815--0.880) & 0.512 & 0.762 & 0.163 & +0.004 (0.594) \\
Stacking + RFECV & KNN ($k=10$) & 13.6 & 0.849 (0.814--0.884) & 0.501 & 0.759 & 0.163 & +0.005 (0.448) \\
\bottomrule

\end{tabular}}
\end{table*}
